% subsection: 階差型 \, $a_{n+1}=a_n+f(n)$

  \subsection{\texorpdfstring{階差型 \, $a_{n+1}=a_n+f(n)$}{階差型}}\label{節：階差型}
  \begin{bookpoint}
    \textbf{変化を足す}\par
    差が添字ごとに変わっても,初項からの変化を足せば元の項に戻る.和の範囲は実際の更新の回数に合わせる.
  \end{bookpoint}
    \bookterm{sequence}{数列}に足す量$f(n)$が番号によって変わっても,順に代入できる.\footnote{$f(n)$は番号$n$の式である.$f(a_n)$のように未知の項の値に依存する場合は,和を書くだけでは計算が済まない.}
    足す量をまとめるため,和の記号を使う.
    \[
      \sum_{k=1}^{n}f(k)=f(1)+f(2)+\cdots+f(n)
    \]
    $k$は和の中で順に動く番号であり,$n$は和の上端である.
    記号の性質は\sectionref*{節：シグマの計算}で確認できる.$n\ge2$で更新を繰り返すと,
    \begin{align*}
      a_n&=a_{n-1}+f(n-1)\\
         &=a_1+f(1)+\cdots+f(n-1)\\
         &=a_1+\sum_{k=1}^{n-1}f(k)
    \end{align*}
    更新は$n-1$回なので,和の上端も$n-1$である.$n=1$の空の和を$0$とすれば,初項も表せる.
    残るのは和の計算である.$f(n)$に応じて,\bookterm{arithmetic}{等差数列}や\bookterm{geometric}{等比数列}の和の公式,\bookterm{partial}{部分分数分解}などを使う.

    \noindent \bookblocklink{example-difference-1}{problem}{answer}{【例題】}\par\nobreak
    \begin{enumerate}[label=(\arabic*),parsep=3pt,format=\bookitemlink{example-difference-1}{problem}{answer}]
      \item $\displaystyle a_1=1,\, a_{n+1}=a_n+n$
      \item $\displaystyle a_1=1,\, a_{n+1}=a_n+2^n$
      \item $\displaystyle a_1=0,\, a_{n+1}=a_n+\frac{1}{n(n+1)}$
      \item $\displaystyle a_1=1,\, a_{n+1}=a_n+(n+1)^2$
    \end{enumerate}

    \noindent\bookblocklink{example-difference-1}{hint}{problem}{【方針】}\par\nobreak
    \begin{enumerate}[label=(\arabic*),parsep=3pt,format=\bookitemlink{example-difference-1}{hint}{problem}]
      \item 差を足し戻し,等差数列の和を使います.
      \item 指数の差も足し戻せます.等比数列の和を使います.
      \item \bookterm{sigma}{シグマ}は\bookterm{partial}{部分分数分解}で処理しましょう.
      \item 初項の$1$も含めると,$1^2+2^2+\cdots+n^2$という平方の和になります.
    \end{enumerate}

    \noindent\bookblocklink{example-difference-1}{answer}{problem}{【解答】}\par\nobreak
    \begin{enumerate}[label=(\arabic*),parsep=3pt,format=\bookitemlink{example-difference-1}{answer}{problem}]
      \item $\displaystyle a_n=1+\frac{n(n-1)}{2}$
      \item $\displaystyle a_n=2^n-1$
      \item $\displaystyle a_n=1-\frac{1}{n}$
      \item $\displaystyle a_n=\frac{1}{6}n(n+1)(2n+1)$
    \end{enumerate}

    \noindent\bookblocklink{example-difference-1}{solution}{problem}{【解法】}\par\nobreak
    \begin{enumerate}[label=(\arabic*),parsep=3pt,format=\bookitemlink{example-difference-1}{solution}{problem}]
      \item $\displaystyle a_1=1,\,a_{n+1}=a_n+n$\\
            \bookterm{difference}{階差}型の\bookterm{general}{一般項}と等差数列の和から,
            \[
              a_n=1+\sum_{k=1}^{n-1}k=1+\frac{n(n-1)}2.
            \]
            和は項数と初項・末項の平均を掛けて求めた.導出は\sectionref*{節：等差数列とその総和の導出}を参照.$n=1$でも初項に一致する.

      \item $\displaystyle a_1=1,\,a_{n+1}=a_n+2^n$\\
            同様に,
            \[
              a_n=1+\sum_{k=1}^{n-1}2^k=1+(2^n-2)=2^n-1.
            \]
            等比数列の和の公式と,倍率をそろえて引く導出は\sectionref*{節：等比数列の総和とその導出}にまとめた.

      \item $\displaystyle a_1=0,\,a_{n+1}=a_n+\frac{1}{n(n+1)}$\\
            \bookterm{partial}{部分分数分解}を用いると,
            \[
              \frac{1}{k(k+1)}
              =\frac{1}{k}-\frac{1}{k+1}
            \]
            と分ければ隣の項が消える.係数の求め方は\sectionref*{節：部分分数分解}を参照.
            \begin{align*}
              a_n
                &=\sum_{k=1}^{n-1}\frac{1}{k(k+1)}\\
                &=\left(1-\frac12\right)+\left(\frac12-\frac13\right)\\
                &\quad+\cdots+\left(\frac1{n-1}-\frac1n\right)\\
                &=1-\frac1n
            \end{align*}

      \item $\displaystyle a_1=1,\,a_{n+1}=a_n+(n+1)^2$\\
            \bookterm{difference}{階差}型の\bookterm{general}{一般項}を用いると,
            \begin{align*}
              a_n
                &=1+\sum_{k=1}^{n-1}(k+1)^2\\
                &=\sum_{j=1}^{n}j^2\\
                &=\frac{1}{6}n(n+1)(2n+1)
            \end{align*}
            最後に平方和の公式を用いた.\booknavref{節：シグマの計算}{公式の一覧}で確認できる.\footnote{公式を差から作る導出は,\booknavref{節：差分から始まる離散的な微積分}{和と差から微積分へ}で扱う.}
    \end{enumerate}
